\begin{proposition}[B.4.4 (Proved): extension of the marked fibration]
Suppose $g:\mathcal Z\to\Delta_s$ is a proper holomorphic map
from a smooth complex fourfold, with reduced normal-crossings
central fibre, and its restriction over $\Delta_s^*$ is
isomorphic to a finite base change of the specified six-sphere
family. The given marked holomorphic map
\[
 f^*:\mathcal Z|_{\Delta_s^*}\longrightarrow\PP^1_t
\]
extends to a meromorphic map $\mathcal Z\dashrightarrow\PP^1_t$.
Its graph gives a proper projective modification
$\mathcal Z'\to\mathcal Z$, an isomorphism over
$\Delta_s^*$, on which the marked map is holomorphic.
No assertion that $\mathcal Z'$ is smooth or semistable is
included in this statement.
\end{proposition}
\begin{proof}
We shrink the disc whenever necessary. Put
$L=K_{\mathcal Z/\Delta_s}^{-2}$. The relative canonical
bundle here is the line bundle
$K_{\mathcal Z}\otimes g^*K_{\Delta_s}^{-1}$; its
restriction to a nonzero fibre is that fibre's canonical bundle.
The precise R1 input is
$K_{X_s}^{-2}\simeq f_s^*\mathcal O_{\PP^1}(1)$,
with connected fibres and no further sections. Therefore
$h^0(X_s,L|_{X_s})=2$, and its complete linear system is
the given map $f_s$, up to a projective change of coordinate.
The three critical values are intrinsically distinguished by
their fibres: multiplicity three, multiplicity four, and reduced
normal crossings, respectively.

Grauert's proper direct-image theorem
\citep[Section 6, Hauptsatz I]{Grauert60push} makes
$\mathcal F=g_*L$ coherent. It is torsion-free: multiplication
by a nonzero base germ is injective on $L$, because the local
semistable equation makes $g$ flat, and hence is injective on
its direct image. A coherent torsion-free sheaf on the smooth
disc is locally free. Its rank is two. On the punctured disc,
constant $h^0=2$ and the analytic base-change theorem checked
in E.3 identify its fibres with the complete spaces of sections.
Choose a frame $(\sigma_0,\sigma_1)$ of $\mathcal F$.
Evaluation gives two global holomorphic sections of $L$ on
$\mathcal Z$, and thus a meromorphic map
$F_0=[\sigma_0:\sigma_1]$ to $\PP^1$.

For completeness its graph is analytic and proper without an
appeal to an unproved extension of a punctured graph. In a local
trivialization of $L$, write the sections as $a,b$. Remove their
common divisorial factor in the regular local ring of the smooth
space $\mathcal Z$. The equation
$v a-u b=0$ in $\mathcal Z\times\PP^1_{[u:v]}$
then has a unique component dominating $\mathcal Z$; it is
the closure of the graph where $(a,b)\ne(0,0)$. These local
components glue. Denote the resulting reduced graph by
$\mathcal G$. The projection $\mathcal G\to\mathcal Z$
is proper and projective, and is an isomorphism over
$\Delta_s^*$. The map
\[
 H:\mathcal G\longrightarrow\PP^1\times\Delta_s,
 \qquad H=(F_0,g),
\]
is proper. Locally $\mathcal G$ is a hypersurface in the smooth
space $\mathcal Z\times\PP^1$. The vanishing of the maximal
Jacobian minors for its equation and the two target coordinates
defines the closed analytic locus $C$ where $H$ is not a
submersion; singular points of $\mathcal G$ are included.

The image $H(C)$ is a closed analytic subset of the target.
One may deduce this precise proper-image assertion from the same
direct-image theorem: the coherent sheaf
$H_*\mathcal O_C$ has support exactly $H(C)$, since the
constant section $1$ is nonzero at every nonempty proper fibre.
Over $\Delta_s^*$ this image consists exactly of three distinct
critical-value sections. It cannot have a two-dimensional
component, since such an analytic component would contain an
open set of the connected two-dimensional target and hence meet
the punctured region in dimension two. Remove any component
contained in $\PP^1\times\{0\}$. The remaining reduced
analytic curve is proper with zero-dimensional fibres over the
disc, and is therefore finite over it.

Each of the three labelled sections on the punctured disc belongs
to a degree-one branch of that finite curve. Indeed their labels
are intrinsic and cannot be permuted by continuation. The
normalization of a degree-$d$ branch over a punctured disc has
cyclic monodromy of length $d$; a single-valued labelled section
therefore has $d=1$. The normalization extends it holomorphically
to a map from the full disc to $\PP^1$. Let these three maps be
$c_1(s),c_2(s),c_0(s)$, in the multiplicity-three,
multiplicity-four and cusp order. They may collide at $s=0$.

For $s\ne0$ the unique projective transformation sending
$(c_1,c_2,c_0)$ to $(0,1,\infty)$ has meromorphic coefficients
at $0$. In a common affine chart its value at $z$ is
\[
 A_s(z)=
 \frac{(z-c_1(s))(c_2(s)-c_0(s))}
      {(z-c_0(s))(c_2(s)-c_1(s))};
\]
homogeneous brackets give the same formula if a value is infinite.
The denominators are not identically zero, since the three values
are distinct on the punctured disc. Thus $A_s\circ F_0$ is
meromorphic on $\mathcal Z$. On every nonzero fibre it agrees
with the actual marked $f_s$, because both are the complete
bi-anticanonical map and have the same three labelled values.
This equality is an equality of maps over the punctured disc,
not just a fibrewise biholomorphism statement.

Finally the graph of $A_s\circ F_0$ is constructed by the same
local two-section argument after clearing powers of the base
coordinate in its meromorphic coefficients. Its projection to
$\mathcal Z$ is the asserted proper projective modification,
and its second projection is the required holomorphic marked map.
\end{proof}

The proposition removes a possible essential projective-coordinate
ambiguity. It does not extend a punctured isomorphism to any of
the auxiliary models of B.3, and does not make a component over
$t=\infty,s=0$ Moishezon. Those would be additional statements.
